\noindent
The map (\ref{equation-sheafy}) is a special case of a base change map.
Namely, suppose that we have a cartesian diagram
$$
\xymatrix{
X' \ar[r]_{g'} \ar[d]_{f'} & X \ar[d]^f \\
Y' \ar[r]^g & Y
}
$$
of quasi-compact and quasi-separated schemes, i.e., a diagram as in
(\ref{equation-base-change}). Assume $f$ and $g$ are {\bf Tor independent}.
Then we can consider the morphism of functors
$D_\QCoh(\mathcal{O}_Y) \to D_\QCoh(\mathcal{O}_{X'})$
given by the composition
\begin{equation}
\label{equation-base-change-map}
L(g')^* \circ a \to
L(g')^* \circ a \circ Rg_* \circ Lg^* \leftarrow
L(g')^* \circ Rg'_* \circ a' \circ Lg^* \to a' \circ Lg^*
\end{equation}
The first arrow comes from the adjunction map $\text{id} \to Rg_* Lg^*$
and the last arrow from the adjunction map $L(g')^*Rg'_* \to \text{id}$.
We need the assumption on Tor independence to invert the arrow
in the middle, see Lemma \ref{lemma-flat-precompose-pus}.
Alternatively, we can think of (\ref{equation-base-change-map}) by
adjointness of $L(g')^*$ and $R(g')_*$ as a natural transformation
$$
a \to a \circ Rg_* \circ Lg^* \leftarrow Rg'_* \circ a' \circ Lg^*
$$
were again the second arrow is invertible. If $M \in D_\QCoh(\mathcal{O}_X)$
and $K \in D_\QCoh(\mathcal{O}_Y)$
then on Yoneda functors this map is given by
\begin{align*}
\Hom_X(M, a(K))
& =
\Hom_Y(Rf_*M, K) \\
& \to
\Hom_Y(Rf_*M, Rg_* Lg^*K) \\
& =
\Hom_{Y'}(Lg^*Rf_*M, Lg^*K) \\
& \leftarrow
\Hom_{Y'}(Rf'_* L(g')^*M, Lg^*K) \\
& =
\Hom_{X'}(L(g')^*M, a'(Lg^*K)) \\
& =
\Hom_X(M, Rg'_*a'(Lg^*K))
\end{align*}
(were the arrow pointing left is invertible by the base
change theorem given in
Derived Categories of Schemes, Lemma \ref{perfect-lemma-compare-base-change})
which makes things a little bit more explicit.